\section{AHSME 1993}
        \begin{enumerate}
        	\item (\textit{AHSME 1993}) For integers $a, b$ and $c$, define $\boxed{a,b,c}$ to mean $a^b-b^c+c^a$. Then $\boxed{1,-1,2}$ equals


 $\text{(A)} \ -4 \qquad \text{(B)} \ -2 \qquad \text{(C)} \ 0 \qquad \text{(D)} \ 2 \qquad \text{(E)} \ 4$

 \textbf{Answer: }D



\textbf{Solution}

Plug in the values for $a,b,c$ and you get $1^{-1} - (-1)^2 + 2^1 \Rightarrow 1-1+2 \Rightarrow \fbox{2}$
$\fbox{D}$



	\item (\textit{AHSME 1993}) In $\triangle ABC$, $\angle A=55^\circ$, $\angle C=75^\circ$, $D$ is on side $\overline{AB}$ and $E$ is on side $\overline{BC}$ If $DB=BE$, then $\angle BED=$
$\text{(A)}\ 50^\circ \qquad \text{(B)}\ 55^\circ \qquad \text{(C)}\ 60^\circ \qquad \text{(D)}\ 65^\circ \qquad \text{(E)}\ 70^\circ$

 \textbf{Answer: }D



\textbf{Solution}

We first consider $\angle CBA$. Because $\angle A = 55$ and $\angle C = 75$, $\angle B = 180 - 55 - 75 = 50$. Then, because $\triangle BED$ is isosceles, we have the equation $2 \angle BED + 50 = 180$. Solving this equation gives us $\angle BED = 65 \rightarrow \fbox{\textbf{(D)}65}$



	\item (\textit{AHSME 1993}) $\frac{15^{30}}{45^{15}} =$
$\text{(A) } \left(\frac{1}{3}\right)^{15}\quad \text{(B) } \left(\frac{1}{3}\right)^{2}\quad \text{(C) } 1\quad \text{(D) } 3^{15}\quad \text{(E) } 5^{15}$

 \textbf{Answer: }E



\textbf{Solution}

First we must convert these to the same bases. We can rewrite $45^{15}$ as $15^{15} \cdot 3^{15}$ Now

 $\frac{15^{30}}{15^{15} \cdot 3^{15}}$ 

 Canceling.... 

 $\frac{15^{15}}{3^{15}} \Rightarrow (\frac{15}{3})^{15} \Rightarrow 5^{15}$
$\fbox{E}$



	\item (\textit{AHSME 1993}) Define the operation "$\circ$" by $x\circ y=4x-3y+xy$, for all real numbers $x$ and $y$. For how many real numbers $y$ does $3\circ y=12$?


 $\text{(A) } 0\quad \text{(B) } 1\quad \text{(C) } 3\quad \text{(D) } 4\quad \text{(E) more than 4}$

 \textbf{Answer: }E



\textbf{Solution}

Note that $3 \circ y = 4 \cdot 3 - 3y + 3y = 12$, so $3 \circ y = 12$ is true for all values of $y$. Thus there are more than four solutions, and the answer is $\fbox{E}$



	\item (\textit{AHSME 1993}) Last year a bicycle cost $160 and a cycling helmet $40. This year the cost of the bicycle increased by $5\%$, and the cost of the helmet increased by $10\%$. The percent increase in the combined cost of the bicycle and the helmet is:


 $\text{(A) } 6\%\quad \text{(B) } 7\%\quad \text{(C) } 7.5\%\quad \text{(D) } 8\%\quad \text{(E) } 15\%$

 \textbf{Answer: }A



\textbf{Solution}

Since the bicycle originally cost $\$160$, the new cost will be $160 \times 5\% = 160 \times \frac{1}{20} = 8$, and $160 + 8 = 168$. The new cost of the helmet will be $40 \times 10\% = 40 \times \frac{1}{10} = 4$, and $40 + 4 = 44$. $160 + 40 = 200$, and $168 + 44 = 212$, so the total percent increase is $\frac{12}{200} = 6\%$, or $\textbf{(A)}$.

 


	\item (\textit{AHSME 1993}) $\sqrt{\frac{8^{10}+4^{10}}{8^4+4^{11}}}=$
$\text{(A) } \sqrt{2}\quad \text{(B) } 16\quad \text{(C) } 32\quad \text{(D) } (12)^{\frac{2}{3}}\quad \text{(E) } 512.5$

 \textbf{Answer: }B



\textbf{Solution 1}

$\sqrt{\frac{ 8^{10}+4^{10} }{8^4 + 4^{11}}} = \sqrt{\frac{2^{30}+2^{20}}{2^{12}+2^{22}}}= \sqrt{\frac{2^{20}(2^{10}+1)}{2^{12}(1+2^{10})}} = \sqrt{2^8} = 2^4$
$\fbox{B}$




\textbf{Solution 2}

$8^{10}+4^{10} = 1074790400$
$8^4 + 4^{11} = 4198400$
$\frac{1074790400}{4198400} = 256$
$\sqrt{256} = 16 = \boxed{B}$



	\item (\textit{AHSME 1993}) The symbol $R_k$ stands for an integer whose base-ten representation is a sequence of $k$ ones. For example, $R_3=111,R_5=11111$, etc. When $R_{24}$ is divided by $R_4$, the quotient $Q=R_{24}/R_4$ is an integer whose base-ten representation is a sequence containing only ones and zeroes. The number of zeros in $Q$ is:


 $\text{(A) } 10\quad \text{(B) } 11\quad \text{(C) } 12\quad \text{(D) } 13\quad \text{(E) } 15$

 \textbf{Answer: }E



\textbf{Solution}

Note $R_n = \sum_{k=0}^{n-1} 10^k = \frac{10^n - 1}{10-1}$.  

 Therefore $\frac{R_{24}}{R_4} = \frac{ 10^{24}-1 }{10^4-1}$.  

 We can recognize this is also the formula for the sum of a geometric series $1+10^4 + (10^4)^2 + \dots + (10^4)^5 = 1+ 10^4 + 10^8 + \dots + 10^{20}$

 Now the 1's place has a 1, but the 10's, 100's and 1,000's place have 0's.  The 10,000's place has a 1, but the $10^5$, $10^6$ and $10^7$ places have 0's.  Between successive 1's in the decimal expansion, there are three 0's, which gives $5\times 3=15$ zeros altogether.

 The answer is $\fbox{E}$



	\item (\textit{AHSME 1993}) Let $C_1$ and $C_2$ be circles of radius 1 that are in the same plane and tangent to each other. How many circles of radius 3 are in this plane and tangent to both $C_1$ and $C_2$?


 $\text{(A) } 2\quad \text{(B) } 4\quad \text{(C) } 5\quad \text{(D) } 6\quad \text{(E) } 8$

 \textbf{Answer: }D



\textbf{Solution}

There are two radius 3 circles to which $C_1$ and $C_2$ are both externally tangent.  One touches the tops of $C_1$ and $C_2$ and extends upward, and the other the other touches the bottoms and extends downward.  There are also two radius 3 circles to which $C_1$ and $C_2$ are both internally tangent, one touching the tops and encircling downward, and the other touching the bottoms and encircling upward.  There are two radius 3 circles passing through the point where $C_1$ and $C_2$ are tangent.  For one $C_1$ is internally tangent and $C_2$ is externally tangent, and for the other $C_1$ is externally tangent and $C_2$ is internally tangent.

 $\fbox{D}$



	\item (\textit{AHSME 1993}) Country $A$ has $c\%$ of the world's population and $d\%$ of the worlds wealth. Country $B$ has $e\%$ of the world's population and $f\%$ of its wealth. Assume that the citizens of $A$ share the wealth of $A$ equally,and assume that those of $B$ share the wealth of $B$ equally. Find the ratio of the wealth of a citizen of $A$ to the wealth of a citizen of $B$.


 $\text{(A) } \frac{cd}{ef}\quad \text{(B) } \frac{ce}{ef}\quad \text{(C) } \frac{cf}{de}\quad \text{(D) } \frac{de}{cf}\quad \text{(E) } \frac{df}{ce}$

 \textbf{Answer: }D



\textbf{Solution}

Let $W$ be the wealth of the world and $P$ be the population of the world.  Hence the wealth of each citizen of $A$ is $w_A = \frac{0.01d W}{0.01cP}=\frac{dW}{cP}$.  Similarly the wealth of each citizen of $B$ is $w_B =\frac{eW}{fP}$.  We divide $\frac{w_A}{w_B} = \frac{de}{cf}$ and see the answer is $\fbox{D}$



	\item (\textit{AHSME 1993}) Let $r$ be the number that results when both the base and the exponent of $a^b$ are tripled, where $a,b>0$. If $r$ equals the product of $a^b$ and $x^b$ where $x>0$, then $x=$
$\text{(A) } 3\quad \text{(B) } 3a^2\quad \text{(C) } 27a^2\quad \text{(D) } 2a^{3b}\quad \text{(E) } 3a^{2b}$

 \textbf{Answer: }C



\textbf{Solution}

We have $r=(3a)^{3b}$

 From this we have the equation $(3a)^{3b}=a^bx^b$

 Raising both sides to the $\frac{1}{b}$ power we get that $27a^3=ax$ or  $x=27a^2$
$\fbox{C}$



	\item (\textit{AHSME 1993}) If $\log_2(\log_2(\log_2(x)))=2$, then how many digits are in the base-ten representation for $x$?


 $\text{(A) } 5\quad \text{(B) } 7\quad \text{(C) } 9\quad \text{(D) } 11\quad \text{(E) } 13$

 \textbf{Answer: }A



\textbf{Solution}

Taking successive exponentials $\log_2(\log_2(x)) = 2^2 = 4$ and $\log_2(x) = 2^4=16$ and $x = 2^{16}$.  Now $2^{10} = 1024 \approx 10^3$ and $2^6 = 64$ so we can approximate $2^{16} \approx 64000$ which has 5 digits.  In general, $2^n$ has approximately $n/3$ digits.

 $\fbox{A}$



	\item (\textit{AHSME 1993}) If $f(2x)=\frac{2}{2+x}$ for all $x>0$, then $2f(x)=$
$\text{(A) } \frac{2}{1+x}\quad \text{(B) } \frac{2}{2+x}\quad \text{(C) } \frac{4}{1+x}\quad \text{(D) } \frac{4}{2+x}\quad \text{(E) } \frac{8}{4+x}$

 \textbf{Answer: }E



\textbf{Solution}

As $f(2x)=\frac{2}{2+x}$, we have that $f(x)=\frac{2}{2+\frac{x}{2}}$. This also means that $2f(x)=\frac{4}{2+\frac{x}{2}}$ which simplifies to $\fbox{E}$. 

 Note: Wait what

 


	\item (\textit{AHSME 1993}) A square of perimeter 20 is inscribed in a square of perimeter 28. What is the greatest distance between a vertex of the inner square and a vertex of the outer square?


 $\text{(A) } \sqrt{58}\quad \text{(B) } \frac{7\sqrt{5}}{2}\quad \text{(C) } 8\quad \text{(D) } \sqrt{65}\quad \text{(E) } 5\sqrt{3}$

 \textbf{Answer: }D



\textbf{Solution}

Assume one of the segments bisected by the inscribed square has length $x$. Thus, the alternate segment has length $7-x$. Applying Pythagorean's Theorem, $x^2+(x-7)^2=5^2$. Simplifying, $(x-3)(x-4)=0$, so $x=3$ or $x=4$ (it does not matter, as rotations produce the same figure). The longest line that can be made forms a right triangle with legs of $4$ and $7$. $\sqrt{4^2+7^2}=\sqrt{65} \rightarrow \boxed{D}$



	\item (\textit{AHSME 1993}) 
                    \begin{figure}[H]
                        \centering
                        \includegraphics[width=0.5\textwidth]{Problem/AHSME_I_1993_14_Problem_Figure_1.png}
                    \end{figure}
                    


 The convex pentagon $ABCDE$ has $\angle{A}=\angle{B}=120^\circ,EA=AB=BC=2$ and $CD=DE=4$. What is the area of ABCDE?


 $\text{(A) } 10\quad \text{(B) } 7\sqrt{3}\quad \text{(C) } 15\quad \text{(D) } 9\sqrt{3}\quad \text{(E) } 12\sqrt{5}$

 \textbf{Answer: }B



\textbf{Solution}


                    \begin{figure}[H]
                        \centering
                        \includegraphics[width=0.5\textwidth]{Solution/AHSME_I_1993_14_Solution_Figure_1.png}
                    \end{figure}
                    

 First, drop perpendiculars from points $A$ and $B$ to segment $EC$.

 Since $\angle EAB = 120^{\circ}$, $\angle EAF = 30^{\circ}$.

 This implies that $\triangle EAF$ is a 30-60-90 Triangle, so $EF = 1$ and $AF = \sqrt3$.

 Similarly, $GB = \sqrt3$ and $GC = 1$.

 Since $FABG$ is a rectangle, $FG = AB = 2$.

 Now, notice that since $EC = EF + FG + GC = 1+2+1=4$, triangle $DEC$ is equilateral.

 Thus, $[ABCDE] = [EABC]+[DCE] = \frac{2+4}{2}(\sqrt3)+\frac{4^2\cdot\sqrt3}{4} = 3\sqrt3+4\sqrt3=\boxed{7\sqrt3 (B)}$

 -AOPS81619

 


	\item (\textit{AHSME 1993}) For how many values of $n$ will an $n$-sided regular polygon have interior angles with integral measures?


 $\text{(A) } 16\quad \text{(B) } 18\quad \text{(C) } 20\quad \text{(D) } 22\quad \text{(E) } 24$

 \textbf{Answer: }D



\textbf{Solution}

Start with the facts that all polygons have their exterior angles sum to 360 and the exterior and interior angles make a linear pair of angles. So our goal is to find the number of divisors of 360 to make both the interior and exterior angles integers. The prime factorization of 360 is $2^3 * 3^2 * 5$. That means the number of divisors is 4*3*2 = 24. But we're not done yet. We cannot have a 1 or 2 sided polygon so we subtract off two bringing us to our final answer of 22 $\fbox{D}$.

 


	\item (\textit{AHSME 1993}) Consider the non-decreasing sequence of positive integers

 \[1,2,2,3,3,3,4,4,4,4,5,5,5,5,5,\cdots\]
in which the $n^{th}$ positive integer appears $n$ times. The remainder when the $1993^{rd}$ term is divided by $5$ is


 $\text{(A) } 0\quad \text{(B) } 1\quad \text{(C) } 2\quad \text{(D) } 3\quad \text{(E) } 4$

 \textbf{Answer: }D



\textbf{Solution}

The sequence of 1's ends at position 1, and the sequence of 2's ends at position 1+2, and the sequence of $n$'s ends at position $1+2+\dots+n$.

 Therefore we want to find the largest integer $n$ that satisfies $\frac{n(n+1)}{2}\geq 1993$.

 By trial and error, the value of $n$ is $63$, and $63 \div 5$ has a remainder of $3$.

 $\fbox{D}$



	\item (\textit{AHSME 1993}) 
                    \begin{figure}[H]
                        \centering
                        \includegraphics[width=0.5\textwidth]{Problem/AHSME_I_1993_17_Problem_Figure_1.png}
                    \end{figure}
                    


 Amy painted a dartboard over a square clock face using the "hour positions" as boundaries.[See figure.] If $t$ is the area of one of the eight triangular regions such as that between 12 o'clock and 1 o'clock, and $q$ is the area of one of the four corner quadrilaterals such as that between 1 o'clock and 2 o'clock, then $\frac{q}{t}=$
$\text{(A) } 2\sqrt{3}-2\quad \text{(B) } \frac{3}{2}\quad \text{(C) } \frac{\sqrt{5}+1}{2}\quad \text{(D) } \sqrt{3}\quad \text{(E) } 2$

 \textbf{Answer: }A



\textbf{Solution}

Assume the length of the side of the square is 2, WLOG.  This means the side of one t section is 1.  As the lines are at clock face positions, each section has a $\frac{360}{12} = 30$ degree angle from the center.  So each section t is a $30-60-90$ triangle with a long leg of 1.  Therefore, the short leg is $\frac{1}{\sqrt3}$.

 
This makes the area of each $t = \frac{1}{2}\cdot b \cdot h = \frac{1}{2}\cdot 1 \cdot \frac{1}{\sqrt3} =  \frac{1}{2\sqrt3}$

 
The total area comprises $4q+8t$, so 
 \[4q+(8\cdot \frac{1}{2\sqrt3}) = 2^2=4\]

 \[4q + \frac{4}{\sqrt3} = 4\]

 \[4q = 4 - \frac{4}{\sqrt3}\]

 \[q = 1 - \frac{1}{\sqrt3} = \frac{\sqrt3 - 1 }{\sqrt3}\]

 \[\frac{q}{t} = \frac{\frac{\sqrt3 - 1 }{\sqrt3}}{\frac{1}{2\sqrt3}} = 2\cdot (\sqrt3 - 1) = \boxed{2\sqrt3-2}\]
$\fbox{A}$



	\item (\textit{AHSME 1993}) Al and Barb start their new jobs on the same day. Al's schedule is 3 work-days followed by 1 rest-day. Barb's schedule is 7 work-days followed by 3 rest-days. On how many of their first 1000 days do both have rest-days on the same day?


 $\text{(A) } 48\quad \text{(B) } 50\quad \text{(C) } 72\quad \text{(D) } 75\quad \text{(E) } 100$

 \textbf{Answer: }E



\textbf{Solution}

We note that Al and Barb's work schedules clearly form a cycle, with Al's work cycle being 4 days long and Barb's work schedule being 10 days long.  Since $\text{LCM}\left(4,10\right)=20$, we know that Al's and Barb's work cycles coincide, and thus collectively repeat, every 20 days.  As a result, we only need to figure out how many coinciding rest days they have in the first 20 days, and multiply that number by $\frac{1000}{20}=50$ to determine the total number of coinciding rest days.  These first 20 days are written out below:

 Day 1:  Al works; Barb works

 Day 2:  Al works; Barb works

 Day 3:  Al works; Barb works

 Day 4:  Al rests; Barb works

 Day 5:  Al works; Barb works

 Day 6:  Al works; Barb works

 Day 7:  Al works; Barb works

 Day 8:  Al rests; Barb rests

 Day 9:  Al works; Barb rests

 Day 10: Al works; Barb rests

 Day 11: Al works; Barb works

 Day 12: Al rests; Barb works

 Day 13: Al works; Barb works

 Day 14: Al works; Barb works

 Day 15: Al works; Barb works

 Day 16: Al rests; Barb works

 Day 17: Al works; Barb works

 Day 18: Al works; Barb rests

 Day 19: Al works; Barb rests

 Day 20: Al rests; Barb rests

 In this 20-day cycle, Al and Barb's rest days coincide twice:  on day 8 and day 20.  As a result (as discussed above), the answer is $2\cdot50=\fbox{100\text{ (E)}}$.    ~cw357

 


	\item (\textit{AHSME 1993}) How many ordered pairs $(m,n)$ of positive integers are solutions to

 \[\frac{4}{m}+\frac{2}{n}=1?\]
$\text{(A) } 1\quad \text{(B) } 2\quad \text{(C) } 3\quad \text{(D) } 4\quad \text{(E) } \text{more than }6$

 \textbf{Answer: }D



\textbf{Solution}

Multiply both sides by $mn$ to clear the denominator. Moving all the terms to the right hand side, the equation becomes $0 = mn-4n-2m$. Adding 8 to both sides allows us to factor the equation as follows: $(m-4)(n-2) = 8$. Since the problem only wants integer pairs $(m,n)$, the pairs are given by the factors of 8, which are 1,2,4, and 8. This results in four pairs, $(5,10), (6,6), (8,4),$ and $(12,3)$, which is answer $\fbox{D}$.

 


	\item (\textit{AHSME 1993}) Consider the equation $10z^2-3iz-k=0$, where $z$ is a complex variable and $i^2=-1$. Which of the following statements is true?


 $\text{(A) For all positive real numbers k, both roots are pure imaginary} \quad\\ \text{(B) For all negative real numbers k, both roots are pure imaginary} \quad\\ \text{(C) For all pure imaginary numbers k, both roots are real and rational} \quad\\ \text{(D) For all pure imaginary numbers k, both roots are real and irrational} \quad\\ \text{(E) For all complex numbers k, neither root is real}$

 \textbf{Answer: }B



\textbf{Solution}

Let $r_1$ and $r_2$ denote the roots of the polynomial.  Then $r_1 + r_2 = \dfrac{3i}{10}$ is pure imaginary, so $r_1$ and $r_2$ have offsetting real parts.  Write $r_1 = a + bi$ and $r_2 = -a + ci$.  

 Now $-k = r_1 r_2 = -a^2 -bc + a(c-b)i$.  In the case that $k$ is real, then $a(c-b)=0$ so either $a=0$ or that $b=c$.  In the first case, the roots are pure imaginary and in the second case we have $k = a^2+b^2$, a positive number.

 We can therefore conclude that if $k$ is real and negative, it must be the first case and the roots are pure imaginary.

 It's possible to rule out the other cases by reasoning through the cases, but this is enough to show that $\fbox{B}$ is true.

 


	\item (\textit{AHSME 1993}) Let $a_1,a_2,\cdots,a_k$ be a finite arithmetic sequence with $a_4 +a_7+a_{10} = 17$ and $a_4+a_5+\cdots+a_{13} +a_{14} = 77$. 


 If $a_k = 13$, then $k =$
$\text{(A) } 16\quad \text{(B) } 18\quad \text{(C) } 20\quad \text{(D) } 22\quad \text{(E) } 24$

 \textbf{Answer: }B



\textbf{Solution}

Note that $a_7-3d=a_4$ and $a_7+3d=a_{10}$ where $d$ is the common difference, so $a_4+a_7+a_{10}=3a_7=17$, or $a_7=\frac{17}{3}$.

 Likewise, we can write every term in the second equation in terms of $a_9$, giving us $11a_9=77\implies a_9=7$.

 Then the common difference is $\frac{2}{3}$. Then $a_k-a_9=13-7=6=9\cdot\frac{2}{3}$.

 This means $a_k$ is $9$ terms after $a_9$, so $k=18\implies\boxed{B}$



	\item (\textit{AHSME 1993}) 
                    \begin{figure}[H]
                        \centering
                        \includegraphics[width=0.5\textwidth]{Problem/AHSME_I_1993_22_Problem_Figure_1.png}
                    \end{figure}
                    


 Twenty cubical blocks are arranged as shown. First, 10 are arranged in a triangular pattern; then a layer of 6, arranged in a triangular pattern, is centered on the 10; then a layer of 3, arranged in a triangular pattern, is centered on the 6; and finally one block is centered on top of the third layer. The blocks in the bottom layer are numbered 1 through 10 in some order. Each block in layers 2,3 and 4 is assigned the number which is the sum of numbers assigned to the three blocks on which it rests. Find the smallest possible number which could be assigned to the top block.


 $\text{(A) } 55\quad \text{(B) } 83\quad \text{(C) } 114\quad \text{(D) } 137\quad \text{(E) } 144$

 \textbf{Answer: }C



\textbf{Solution}

The value assigned at the top is the weighted sum of the values in the cubes of a given level.  The weight for a given cube is the sum of the weights of the blocks which touch it above, since the number in a cube will get incorporated into the overall sum via each of those blocks, according to the weight of each block.  Thus, if we think about how the weights propagate down the block pyramid, its sort of like a three dimensional Pascal's triangle.

 The first layer is $1$

 The second layer is 

 $| \, \, 1 \\ |1 \, 1$

 The third layer is 

 $| \,\,\,\, 1 \\| \,\, 2 \, 2 \\| 1 \,2 \, 1$

 The fourth layer is 

 $| \,\,\,\,\,\, 1 \\| \,\,\,\, 3 \, 3 \\| \,\, 3 \,6 \, 3 \\| 1\, 3\, 3\, 1$

 The top level sum is minimized if we associate the block with weight 6 with the smallest number 1, the blocks with weight 1 with the largest numbers 8, 9 and 10, and the rest of the blocks with weight 3 with the rest of the numbers.  The sum is $6\cdot 1 + 3 \cdot (2+3+4+5+6+7 ) + 1\cdot (8+9+10) = 114$ so the answer is  $\fbox{C}$



	\item (\textit{AHSME 1993}) 
                    \begin{figure}[H]
                        \centering
                        \includegraphics[width=0.5\textwidth]{Problem/AHSME_I_1993_23_Problem_Figure_1.png}
                    \end{figure}
                    


 Points $A,B,C$ and $D$ are on a circle of diameter $1$, and $X$ is on diameter $\overline{AD}.$


 If $BX=CX$ and $3\angle{BAC}=\angle{BXC}=36^\circ$, then $AX=$
$\text{(A) } \cos(6^\circ)\cos(12^\circ)\sec(18^\circ)\quad\\ \text{(B) } \cos(6^\circ)\sin(12^\circ)\csc(18^\circ)\quad\\ \text{(C) } \cos(6^\circ)\sin(12^\circ)\sec(18^\circ)\quad\\ \text{(D) } \sin(6^\circ)\sin(12^\circ)\csc(18^\circ)\quad\\ \text{(E) } \sin(6^\circ)\sin(12^\circ)\sec(18^\circ)$

 \textbf{Answer: }B



\textbf{Solution}

We have all the angles we need, but most obviously, we see that right angle in triangle $ABD$.

 Note also that angle $BAD$ is 6 degrees, so length $AB = cos(6)$ because the diameter, $AD$, is 1.

 Now, we can concentrate on triangle $ABX$ (after all, now we can decipher all angles easily and use Law of Sines).

 We get:

 $\frac{AB}{\sin(\angle{AXB})} =\frac{AX}{\sin(\angle{ABX})}$

 That's equal to

 $\frac{\cos(6)}{\sin(180-18)} =\frac{AX}{\sin(12)}$

 Therefore, our answer is equal to:$\fbox{B}$

 Note that $\sin(162) = \sin(18)$, and don't accidentally put $\text{C}$ because you thought $\frac{1}{\sin}$ was $\sec$!

 


	\item (\textit{AHSME 1993}) A box contains $3$ shiny pennies and $4$ dull pennies. One by one, pennies are drawn at random from the box and not replaced. If the probability is $a/b$ that it will take more than four draws until the third shiny penny appears and $a/b$ is in lowest terms, then $a+b=$
$\text{(A) } 11\quad \text{(B) } 20\quad \text{(C) } 35\quad \text{(D) } 58\quad \text{(E) } 66$

 \textbf{Answer: }E



\textbf{Solution}

First let’s try to find the number of possible unique combinations. I’ll denote shiny coins as 1 and dull coins as 0. 

 Now, each configuration can be represented by a string of 1s and 0s e.g. 0110100. Notice that a combination can be uniquely determined solely by the placement of their 0s OR 1s e.g. 1 - - 1 1 - - where the dashes can be replaced by 0s. This makes the number of unique combinations 7 choose 3 (if you’re counting w.r.t. shiny coins) OR 7 choose 4 (w.r.t dull coins). Both are equal to 35.

 Next, observe that, for the event that the third shiny coin is not within your first 4 picks, it has to be within the last three numbers. You can think of this as placing the seven coins in a vertical stack in the box and shuffling that stack randomly. Then, to pick, you extract the first coin on the top and keep repeating. It has the same effect.

 The sequence can have 1 shiny coin in the last 3 digits (Case 1), 2 shiny coins in the last 3 digits (Case 2) or 3 shiny coin in the last three digits (Case 3).

 Case 1:Let’s start with the first case of one shiny coin in its last 3 digits.

 Example: 0110100

 The first four numbers has 4 spaces and 2 shiny coins therefore the number of combinations is 4 choose 2 = 6. The last 3 digits has 3 combinations for the same reason.So, probability for Case 1 to occur is:$\dfrac{6*3}{35}=\dfrac{18}{35}$

 Case 2:Using the fact that the combinations are uniquely determined by an order of 0s or 1s and you can just fill the rest, in you can ascertain:$\underbracket{0100}_{\text{4 combinations}}\, \underbracket{110}_{\text{3 combinations}}$

 So, P(Case 2)=$12/35$

 Case 3:Trivially, it is 1.P(Case 3)=$1/35$

 Adding all these probabilites together gives you the probability that the third shiny coin will not appear in your first 4 draws:$\dfrac{18}{35}+\dfrac{12}{35}+\dfrac{1}{35}=\dfrac{31}{35}$
$\dfrac{a}{b}=\dfrac{31}{35}$

 Since the fraction is irreducible:

 $a=31$,$b=35$

 \[a+b=66\]

 The answer is E.

 


 



\textbf{Solution 2}

Using complementary probability, we can reduce the problem into two cases- the third shiny penny is drawn on the third draw, or on the fourth draw. For the first case, there is only one way to have the shiny pennies as the first three coins drawn, out of a possible 35 drawings (7 choose 3). For the second case, the third shiny penny has to be the fourth penny drawn, which leaves three possible orderings for the first three coins drawn (NS S S, S NS S, S S NS), out of 35 (7 choose 4). Adding these two probabilities together gives $\dfrac{4}{35}$, and subtracting this from one yields $\dfrac{31}{35}$, which makes a 31 and b 35, which sum to $\fbox{E}$.

 


	\item (\textit{AHSME 1993}) 
                    \begin{figure}[H]
                        \centering
                        \includegraphics[width=0.5\textwidth]{Problem/AHSME_I_1993_25_Problem_Figure_1.png}
                    \end{figure}
                    


 Let $S$ be the set of points on the rays forming the sides of a $120^{\circ}$ angle, and let $P$ be a fixed point inside the angle 
on the angle bisector. Consider all distinct equilateral triangles $PQR$ with $Q$ and $R$ in $S$. 
(Points $Q$ and $R$ may be on the same ray, and switching the names of $Q$ and $R$ does not create a distinct triangle.) 
There are


 $\text{(A) exactly 2 such triangles} \quad\\ \text{(B) exactly 3 such triangles} \quad\\ \text{(C) exactly 7 such triangles} \quad\\ \text{(D) exactly 15 such triangles} \quad\\ \text{(E) more than 15 such triangles}$

 \textbf{Answer: }E



\textbf{Solution}

$\fbox{E}$

 Take the "obvious" equilateral triangle $OAP$, where $O$ is the vertex, $A$ is on the upper ray, and $P$ is our central point.  Slide $A$ down on the top ray to point $A'$, and slide $O$ down an equal distance on the bottom ray to point $O'$.

 Now observe $\triangle AA'P$ and $\triangle OO'P$.  We have $m\angle A = 60^\circ$ and $m \angle O = 60^\circ$, therefore $\angle A \cong \angle O$.  By our construction of moving the points the same distance, we have $AA' = OO'$.  Also, $AP = OP$ by the original equilateral triangle.  Therefore, by SAS congruence, $\triangle AA'P \cong \triangle OO'P$.  

 Now, look at $\triangle A'PO'$.  We have $PA' = PO'$ from the above congruence.  We also have the included angle $\angle A'PO'$ is $60^\circ$.  To prove that, start with the $60^\circ$ angle $APO$, subtract the angle $APA'$, and add the congruent angle $OPO'$.

 Since $\triangle A'PO'$ is an isosceles triangle with vertex of $60^\circ$, it is equilateral.

 


	\item (\textit{AHSME 1993}) Find the largest positive value attained by the function

 \[f(x)=\sqrt{8x-x^2}-\sqrt{14x-x^2-48} ,\quad x  \text{  a real number}\]
$\text{(A) } \sqrt{7}-1\quad \text{(B) } 3\quad \text{(C) } 2\sqrt{3}\quad \text{(D) } 4\quad \text{(E) } \sqrt{55}-\sqrt{5}$

 \textbf{Answer: }C



\textbf{Solution}

We can rewrite the function as $f(x) = \sqrt{x (8 - x)} - \sqrt{(x - 6) (8 - x)}$ and then factor it to get $f(x) = \sqrt{8 - x} \left(\sqrt{x} - \sqrt{x - 6}\right)$. From the expressions under the square roots, it is clear that $f(x)$ is only defined on the interval $[6, 8]$.

 The $\sqrt{8 - x}$ factor is decreasing on the interval. The behavior of the $\sqrt{x} - \sqrt{x - 6}$ factor is not immediately clear. But rationalizing the numerator, we find that $\sqrt{x} - \sqrt{x - 6} = \frac{6}{\sqrt{x} + \sqrt{x - 6}}$, which is monotonically decreasing. Since both factors are always positive, $f(x)$ is also positive. Therefore, $f(x)$ is decreasing on $[6, 8]$, and the maximum value occurs at $x = 6$. Plugging in, we find that the maximum value is $\boxed{\text{(C) } 2\sqrt{3}}$.

 



\textbf{Solution 2}

Note the form of the function is $f(x) = \sqrt{ p(x)} - \sqrt{q(x)}$ where $p(x)$ and $q(x)$ each describe a parabola.  Factoring we have $p(x) = x(8-x)$ and $q(x) = (x-6)(8-x)$.  

 The first term $\sqrt{p(x)}$ is defined only when $p(x)\geq 0$ which is the interval $[0,8]$ and the second term $\sqrt{q(x)}$ is defined only when $q(x)\geq 0$ which is on the interval $[6,8]$, so the domain of  $f(x)$ is $[0,8] \cap [6,8] = [6,8]$.  

 Now $p(x)$ peaks at the midpoint of its roots at $x=4$ and it decreases to 0 at $x=8$.  Thus, $p(x)$ is decreasing over the entire domain of $f(x)$ and it   obtains its maximum value at the left boundary $x=6$ and $\sqrt{p(x)}$ does as well.  On the other hand $q(x)$ obtains its minimum value of $q(x)=0$ at the left boundary $x=6$ and $\sqrt{q(x)}$ does as well.  Therefore $\sqrt{p(x)}-\sqrt{q(x)}$ is maximized at $x=6$.  (If this seems a little unmotivated, a quick sketch of the two parabolic-like curves makes it clear where the distance between them is greatest). 

 The value at $x=6$ is $\sqrt{ 6\cdot 2 } = 2\sqrt{3}$ and the answer is $\fbox{C}$.

 


	\item (\textit{AHSME 1993}) 
                    \begin{figure}[H]
                        \centering
                        \includegraphics[width=0.5\textwidth]{Problem/AHSME_I_1993_27_Problem_Figure_1.png}
                    \end{figure}
                    
The sides of $\triangle ABC$ have lengths $6,8,$ and $10$. A circle with center $P$ and radius $1$ rolls around the inside of $\triangle ABC$, always remaining tangent to at least one side of the triangle. When $P$ first returns to its original position, through what distance has $P$ traveled?


 $\text{(A) } 10\quad \text{(B) } 12\quad \text{(C) } 14\quad \text{(D) } 15\quad \text{(E) } 17$

 \textbf{Answer: }B



\textbf{Solution}


                    \begin{figure}[H]
                        \centering
                        \includegraphics[width=0.5\textwidth]{Solution/AHSME_I_1993_27_Solution_Figure_1.png}
                    \end{figure}
                    

 Start by considering the triangle traced by $P$ as the circle moves around the triangle. It turns out this triangle is similar to the $6-8-10$ triangle (Proof: Realize that the slope of the line made while the circle is on $AC$ is the same as line $AC$ and that it makes a right angle when the circle switches from being on $AB$ to $BC$). Then, drop the perpendiculars as shown.

 Since the smaller triangle is also a $6-8-10 = 3-4-5$ triangle, we can label the sides $EF,$ $CE,$ and $DF$ as $3x, 4x,$ and $5x$ respectively. Now, it is clear that $GB = DE + 1 = 4x + 1$, so $AH = AG = 8 - GB = 7 - 4x$ since $AH$ and $AG$ are both tangent to the circle P at some point. We can apply the same logic to the other side as well to get $CI = 5 - 3x$. Finally, since we have $HI = DF = 5x$, we have $AC = 10 = (7 - 4x) + (5x) + (5 - 3x) = 12 - 2x$, so $x = 1$ and $3x + 4x + 5x = \fbox{(B) 12}$

 -Solution by Someonenumber011

 


	\item (\textit{AHSME 1993}) How many triangles with positive area are there whose vertices are points in the $xy$-plane whose coordinates are integers $(x,y)$ satisfying $1\le x\le 4$ and $1\le y\le 4$?


 $\text{(A) } 496\quad \text{(B) } 500\quad \text{(C) } 512\quad \text{(D) } 516\quad \text{(E) } 560$

 \textbf{Answer: }D



\textbf{Solution}

The vertices of the triangles are limited to a $4\times4$ grid, with $16$ points total. Every triangle is determined by $3$ points chosen from these $16$ for a total of $\binom{16}{3}=560$. However, triangles formed by collinear points do not have positive area. For each column or row, there are $\binom{4}{3}=4$ such degenerate triangles. There are $8$ total columns and rows, contributing $32$ invalid triangles. There are also $4$ for both of the diagonals and $1$ for each of the $4$ shorter diagonals. There are a total of $32+8+4=44$ invalid triangles counted in the $560$, so the answer is $560-44=516$ or answer choice $\fbox{D}$.

 


	\item (\textit{AHSME 1993}) Which of the following could NOT be the lengths of the external diagonals of a right regular prism [a "box"]? (An $\textit{external diagonal}$ is a diagonal of one of the rectangular faces of the box.)


 $\text{(A) }\{4,5,6\} \quad \text{(B) } \{4,5,7\} \quad \text{(C) } \{4,6,7\} \quad \text{(D) } \{5,6,7\} \quad \text{(E) } \{5,7,8\}$

 \textbf{Answer: }B



\textbf{Solution}

Let $a,$ $b,$ and $c$ be the side lengths of the rectangular prism. By Pythagoras, the lengths of the external diagonals are $\sqrt{a^2 + b^2},$ $\sqrt{b^2 + c^2},$ and $\sqrt{a^2 + c^2}.$ If we square each of these to obtain $a^2 + b^2,$ $b^2 + c^2,$ and $a^2 + c^2,$ we observe that since each of $a,$ $b,$ and $c$ are positive, then the sum of any two of the squared diagonal lengths must be larger than the square of the third diagonal length. For example, $(a^2 + b^2) + (b^2 + c^2) = (a^2 + c^2) + 2b^2 > a^2 + c^2$ because $2b^2 > 0.$

 Thus, we test each answer choice to see if the sum of the squares of the two smaller numbers is larger than the square of the largest number. Looking at choice (B), we see that $4^2 + 5^2 = 41 < 7^2 = 49,$ so the answer is $\boxed{\text{(B) \{4, 5, 7\}}}.$

 --goldentail141

 



\textbf{Solution 2}

Let $a,$ $b,$ and $c$ be the side lengths of the rectangular prism. We then have:$a^2+b^2=i^2,$$b^2+c^2=j^2,$$c^2+a^2=k^2$where $\{ i,j,k \}$ are the given diagonal lengths. WLOG, let $i \leq j \leq k.$ It suffices to check if $i^2+j^2 \geq k.$ We see that for $\boxed{B}$, $4^2+5^2 = 16 + 25 = 41 < 7^2 = 49$, so this case is impossible.

 


	\item (\textit{AHSME 1993}) Given $0\le x_0<1$, let 

 \[x_n=\left\{ \begin{array}{ll} 2x_{n-1} &\text{ if }2x_{n-1}<1 \\ 2x_{n-1}-1 &\text{ if }2x_{n-1}\ge 1 \end{array}\right.\]
for all integers $n>0$. For how many $x_0$ is it true that $x_0=x_5$?


 $\text{(A) 0} \quad \text{(B) 1} \quad \text{(C) 5} \quad \text{(D) 31} \quad \text{(E) }\infty$

 \textbf{Answer: }D



\textbf{Solution}

We are going to look at this problem in binary. 

 $x_0 = (0.a_1 a_2 \cdots )_2$
$2x_0 = (a_1.a_2 a_3 \cdots)_2$

 If $2x_0 < 1$, then $x_0 < \frac{1}{2}$ which means that $a_1 = 0$ and so $x_1 = (.a_2 a_3 a_4 \cdots)_2$

 If $2x_0 \geq 1$ then $x \geq \frac{1}{2}$ which means that $x_1 = 2x_0 - 1 = (.a_2 a_3 a_4 \cdots)_2$.

 Using the same logic, we notice that this sequence cycles and that since $x_0 = x_5$ we notice that $a_n = a_{n+5}$. 

 We have $2$ possibilities for each of $a_1$ to $a_5$ but we can't have $a_1 = a_2 = a_3 = a_4 = a_5 = 1$ so we have $2^5 - 1 = \boxed{(D)31}$

 -mathman523

 



\textbf{Solution 2}

$x_5 = \left\{ 2^5x_0 \right\}$
$\therefore x_0 = \left\{ 2^5x_0 \right\}$
$Suppose \  x_0 = 2^5 x_0 - k \ (k \geq 0)$
$31x_0 = k$
$x_0 = \frac{k}{31}$
$\therefore k = 0, 1, 2,..., 30$

 select (D)

 


\end{enumerate}